\documentclass[10pt,a4paper]{amsart}
\usepackage[margin=19mm]{geometry}
\usepackage{amsmath,amssymb,mathtools}
\usepackage[scr=rsfs,cal=euler]{mathalfa}
\usepackage{xfrac}
\usepackage[unicode,hidelinks,bookmarks=false]{hyperref}
\newcommand{\ie}{i.e.\ }
\newcommand{\cf}{cf.\ }
\newcommand{\resp}{respectively\ }
\newcommand{\Subsection}{\subsection}
\providecommand{\nobreakdash}{\nobreak\hbox{-}}
\setlength{\emergencystretch}{2em}
\multlinegap0pt
\usepackage{amssymb}
\usepackage{csquotes}
\newtheorem{lemma}{Lemma}
\newtheorem{theorem}{Theorem}
\title{Cryptanalysis of the Legendre Pseudorandom Function over Extension Fields}
\author{Daksh Pandey}
\date{Source: arXiv:2604.04833v3 (CC BY 4.0)}
\begin{document}
\maketitle

\noindent\emph{Source and licence.} Cryptanalysis of the Legendre Pseudorandom Function over Extension Fields. Version arXiv:2604.04833v3 (CC BY 4.0), distributed by arXiv under the Creative Commons Attribution 4.0 International licence, http://creativecommons.org/licenses/by/4.0/, as recorded in the arXiv licence metadata for this exact version and shown on the abstract page https://arxiv.org/abs/2604.04833v3. Full licence notice: \texttt{licenses/arxiv-2604.04833-CC-BY-4.0.txt}.\par\smallskip

\noindent\textit{Editorial scope.} Counted unit: the article's theorem thm:active (source0.tex lines 209--211). The article's complete original proof of it is delivered with it (source0.tex lines 213--227, one \texttt{proof} environment of 15 lines). No mathematics is written, repaired or re-derived: the statement and the proof are the article's own lines, in the article's own notation, and no retained line is substituted or re-flowed.  Retained uncounted as the source-local context the proof needs: lines 190--192.  Nothing was omitted from any retained span, so the package carries no omission record.  No retained line contains a citation, so the wrapper carries no bibliography and no bibliography entry.  No retained line contains a figure environment, an image-inclusion command or drawing code, so no image is delivered and the package declares no asset dependency.  Editorial formatting only, in addition: an AMS \texttt{amsart} preamble that restates the article's own declarations and macros used by the retained lines, namely the environments \texttt{lemma}, \texttt{theorem} on separate counters, together with the standard packages those lines need.  The retained environments print in this document as Lemma 1, Theorem 1 in order of appearance, whereas the article prints the same items with the numbering of its own full document.  The complete native source of the article as distributed in the arXiv e-print for this exact version is bundled unchanged as \texttt{sources/197913/source0.tex}.
\begin{lemma}[Factoring Identity of the Legendre PRF] \label{lemma:factor}
For a geometric query sequence $g^i$, the Legendre PRF evaluation factors into a secret-dependent constant and a universal, keyless reference sequence shifted by a discrete logarithm $j$.
\end{lemma}

\begin{theorem}[Active Key Recovery] \label{thm:active}
Given $M$ chosen geometric queries, an adversary can extract the secret polynomial $K \in \mathbb{F}_{p^r}$ in purely $\mathcal{O}(p^r / M)$ operations.
\end{theorem}

\begin{proof}
The adversary intercepts the geometric keystream and populates a hash table $T$ mapping $U$-bit contiguous windows to their starting index $i$. To mitigate spurious collisions, $U$ is chosen such that $2^U \gg p^r$.

The adversary defines the universal reference function $L_0(m) = L(g^m + 1)$. By randomly sampling exponents $m \in \{1, \dots, p^r-1\}$, the adversary generates reference windows of length $U$. By Lemma \ref{lemma:factor}, the target sequence is a strict shift of the reference sequence, potentially inverted if $L(K) = -1$. Therefore, the adversary queries $T$ for both the generated reference window and its bitwise complement.

Because $T$ contains $M - U \approx M$ targets, a collision is guaranteed in expected $p^r/M$ guesses. Let the collision match target index $i$ and reference index $m$. By the factoring identity, the exponent relation is:
\begin{equation*}
    m \equiv i + j \pmod{p^r - 1} \implies j \equiv m - i \pmod{p^r - 1}
\end{equation*}
The secret polynomial $K$ is then extracted via the field inversion:
\begin{equation*}
    K \equiv (g^j)^{-1} \pmod{I(x), p}
\end{equation*}
This process requires exactly $\mathcal{O}(p^r/M)$ Legendre symbol evaluations, effectively reducing the active security of the single-degree PRF in $\mathbb{F}_{p^r}$ to its theoretical lower bound.
\end{proof}

\end{document}
